\Needspace{30\baselineskip}
\section{সংক্ষিপ্ত প্রশ্ন ও সৃজনশীল অনুশীলনী}
\label{sec:ch06-exercises}
 নিচের প্রশ্নগুলো অধ্যায়ের ধারণা প্রয়োগের জন্য প্রস্তুত করা হয়েছে। প্রতিটি উদ্দীপকে দেওয়া টেবিল তার নিজস্ব নমুনা; একই নামের অন্য টেবিলের মান ব্যবহার করবে না। কোনো প্রশ্নে অধ্যায়ের আগের টেবিল ব্যবহার করতে বলা হলে শুধু সেই নির্দিষ্ট টেবিলটি নেবে।

\subsection{ব্যাখ্যা ও তুলনামূলক প্রশ্ন}
\begin{multicols}{2}
\small
\begin{enumerate}[itemsep=0pt,topsep=1pt]
\item একটি শ্রেণির গড় নম্বর কীভাবে আরেকটি বিশ্লেষণের ডেটা হতে পারে?
\item একই ফোল্ডারে থাকা সব ফাইলকে একটি নির্দিষ্ট ডেটাবেজ বলা যায় না কেন?
\item ফিল্ডের নাম, ফিল্ডের মান ও রেকর্ডের পার্থক্য উদাহরণ দিয়ে বোঝাও।
\item GPA-এর Data Type নির্ধারণের পাশাপাশি Validation Rule কেন প্রয়োজন, ব্যাখ্যা করো।
\item পাঁচ কলামের টেবিলে তিন রেকর্ড যোগ ও এক কলাম বাদ দিলে কোন পরিমাপ কীভাবে বদলায়?
\item DBMS ও DBA-এর দায়িত্ব আলাদা করে ব্যাখ্যা করো।
\item একই নামের দুজনের ক্ষেত্রে Name-কে প্রাইমারি কী করলে সমস্যা কী?
\item একাধিক ফিল্ড পৃথকভাবে অনন্য হলে প্রাইমারি কী নির্বাচনে কী বিবেচনা করবে?
\item একটি প্রাইমারি কী একাধিক ফিল্ড নিয়ে গঠিত হতে পারে কেন?
\item ফরেন কী-এর পুনরাবৃত্তি কীভাবে 1:M সম্পর্ককে প্রকাশ করে?
\item ফিল্ডের নাম আলাদা হলেও সম্পর্ক সম্ভব, কিন্তু অসামঞ্জস্যপূর্ণ টাইপ সমস্যা করে কেন?
\item পরীক্ষায় অনুপস্থিত শিক্ষার্থীর নম্বর এবং শূন্য পাওয়া শিক্ষার্থীর নম্বর কীভাবে আলাদা করবে?
\item Format, Input Mask ও Validation Rule-এর কাজ আলাদা করে বোঝাও।
\item AutoNumber-এর সর্বোচ্চ মানকে মোট রেকর্ডসংখ্যা ধরা যায় না কেন?
\item SELECT-এর expression এবং ALTER TABLE দিয়ে নতুন ফিল্ড যোগের পার্থক্য কী?
\item WHERE ও HAVING কোন পর্যায়ের শর্ত নির্ধারণ করে?
\item ইনডেক্স অতিরিক্ত জায়গা ও লেখার ব্যয় বাড়ায় কেন?
\item NoSQL একটি ভাষা নয় কেন? Big Data-এর সঙ্গে এর পার্থক্য কী?
\item ERP-তে একই অর্ডারের তথ্য বিভিন্ন বিভাগ কীভাবে ব্যবহার করতে পারে?
\item এনক্রিপশন, permission ও backup কেন পরস্পরের পূর্ণ বিকল্প নয়?
\end{enumerate}
\end{multicols}

\subsection{সারণি ও SQL-এর ফল বিশ্লেষণ}
নিচের Product টেবিল ব্যবহার করো। UnitPrice-এর সব মান একই মুদ্রায় এবং Stock পূর্ণসংখ্যায় দেওয়া হয়েছে।
\begin{flushleft}\small
\begin{tabular}{|l|l|l|l|}\hline
ProductID & ProductName & UnitPrice & Stock\\\hline
11 & Notebook & 80 & 30\\\hline
12 & Pen & 15 & 100\\\hline
13 & Folder & 80 & 20\\\hline
14 & Marker & 45 & 0\\\hline
\end{tabular}
\end{flushleft}
\begin{enumerate}
\item টেবিলের ফিল্ড, রেকর্ড ও ডেটা-ঘরের সংখ্যা নির্ণয় করো। UnitPrice একা প্রাইমারি কী হবে না কেন?
\item \begin{minipage}[t]{\linewidth}নিচের নির্দেশগুলোর ফল নির্ণয় করো। ফলের কলামের নাম ও নির্বাচিত রেকর্ড আলাদা করে লেখো।
\begin{verbatim}
SELECT ProductID,ProductName FROM Product
WHERE UnitPrice>=45 AND Stock>0;
SELECT ProductName,UnitPrice*Stock AS StockValue
FROM Product WHERE ProductID IN (11,13);
SELECT ProductID FROM Product ORDER BY UnitPrice DESC,ProductID ASC;
SELECT COUNT(*),SUM(Stock),MAX(UnitPrice) FROM Product;
\end{verbatim}
\end{minipage}
\item Stock শূন্য এমন পণ্যের নাম দেখাও এবং নতুন Stock মান 10 করতে একটি UPDATE লেখো। WHERE ছাড়া চালালে কী হবে?
\item SupplierCode নামে Text ফিল্ড যোগের SQLite নির্দেশ লেখো। এতে কোনো নতুন পণ্যের রেকর্ড যোগ হবে কি?
\item ProductID 12-এর রেকর্ড মুছে ফেলা ও Product টেবিল বাদ দেওয়ার নির্দেশ আলাদা করে লেখো।
\end{enumerate}


\subsection{ত্রুটি শনাক্ত ও সংশোধন}
নিচের প্রতিটি ক্ষেত্রে সমস্যা চিহ্নিত করে সঠিক রূপ ও কারণ লেখো।
\begin{multicols}{2}
\small
\begin{enumerate}
\item অজানা GPA খুঁজতে \texttt{WHERE GPA=NULL} লেখা হয়েছে।
\item শুধু রোল 101 মুছতে \texttt{DELETE Student WHERE Roll=101;} লেখা হয়েছে।
\item নাম বদলাতে \texttt{ALTER TABLE Student SET Name='Rupa';} লেখা হয়েছে।
\item Text মোবাইল নম্বর 01710000000-কে হিসাবের সুবিধার যুক্তিতে Number করা হয়েছে।
\item Enrolment-এ অনেক শিক্ষার্থী বহু কোর্স নেয়, কিন্তু শুধু StudentID-কে unique করা হয়েছে।
\item নতুন কর্মীর পরিচিতি AutoNumber; অন্য টেবিলের EmployeeID-ও AutoNumber করে সম্পর্কিত মান মেলার আশা করা হয়েছে।
\item SQLite-এ নতুন ডেটাবেজ ফাইলের পরিবর্তে \texttt{CREATE DATABASE School;} চালানো হয়েছে।
\item একটি ছোট নমুনায় সব নাম আলাদা দেখে Name-কে চিরকাল অনন্য ঘোষণা করা হয়েছে।
\end{enumerate}
\end{multicols}


\subsection{সৃজনশীল প্রশ্ন}
\textbf{১. গ্রন্থাগারের ডেটা নকশা।} একটি কলেজে Book(BookID, Title, AuthorName) এবং Issue(IssueID, BookID, StudentID, IssueDate, ReturnDate) রাখা হয়। একটি বই ভিন্ন সময়ে বিভিন্ন শিক্ষার্থী নিতে পারে। ফেরত না দেওয়া বইয়ের ReturnDate এখনও জানা নেই।
\begin{multicols}{2}
\begin{enumerate}
\item[ক.] \begin{minipage}[t]{\linewidth}Field কী?
\end{minipage}
\item[খ.] \begin{minipage}[t]{\linewidth}Title-কে বইয়ের অনন্য পরিচিতি ধরা সবসময় যুক্তিযুক্ত নয় কেন?
\end{minipage}
\item[গ.] \begin{minipage}[t]{\linewidth}টেবিল দুটির কী, সম্পর্ক এবং IssueDate ও ReturnDate-এর উপযুক্ত Access টাইপ নির্ধারণ করো।
\end{minipage}
\item[ঘ.] \begin{minipage}[t]{\linewidth}ফেরত না দেওয়া বইয়ের তথ্য নির্ভরযোগ্যভাবে খোঁজার ব্যবস্থা বিশ্লেষণ করো।
\end{minipage}
\end{enumerate}
\end{multicols}

\textbf{২. শ্রেণি ও রোল।} একটি প্রতিষ্ঠানে XI ও XII উভয় শ্রেণিতে Roll 1 থেকে শুরু হয়। একই শিক্ষাবর্ষে প্রতি শ্রেণির একটি শাখা আছে। Student(ClassName, Roll, Name, Mobile, GPA)-তে বর্তমান শিক্ষাবর্ষের তথ্য রাখা হয়। ভবিষ্যতে একাধিক শিক্ষাবর্ষের তথ্য রাখার পরিকল্পনা আছে।
\begin{multicols}{2}
\begin{enumerate}
\item[ক.] \begin{minipage}[t]{\linewidth}কম্পোজিট কী কী?
\end{minipage}
\item[খ.] \begin{minipage}[t]{\linewidth}Mobile দেখতে সংখ্যা হলেও Text উপযুক্ত কেন?
\end{minipage}
\item[গ.] \begin{minipage}[t]{\linewidth}বর্তমান কাঠামোয় উপযুক্ত অনন্য কী ও GPA-এর যাচাইয়ের নিয়ম নির্ধারণ করো।
\end{minipage}
\item[ঘ.] \begin{minipage}[t]{\linewidth}ভবিষ্যৎ শিক্ষাবর্ষ যোগ হলে বর্তমান কী যথেষ্ট থাকবে কি না বিশ্লেষণ করো।
\end{minipage}
\end{enumerate}
\end{multicols}

\textbf{৩. কোর্সে ভর্তি।} TrainingCentre-এ Learner(LearnerID, Name), Course(CourseID, CourseName) এবং Enrolment(LearnerID, CourseID, EnrolDate) আছে। একই শিক্ষার্থী একই কোর্সে একবার ভর্তি হতে পারে, কিন্তু ভিন্ন কোর্সে ভর্তি হওয়ার সুযোগ আছে।
\begin{multicols}{2}
\begin{enumerate}
\item[ক.] \begin{minipage}[t]{\linewidth}জাংশন টেবিল কী?
\end{minipage}
\item[খ.] \begin{minipage}[t]{\linewidth}শুধু LearnerID-কে Enrolment-এর প্রাইমারি কী করলে সমস্যা কী?
\end{minipage}
\item[গ.] \begin{minipage}[t]{\linewidth}তিন টেবিলের কী ও সংযোগচিত্র তৈরি করো।
\end{minipage}
\item[ঘ.] \begin{minipage}[t]{\linewidth}অজানা CourseID গ্রহণ করলে কী সমস্যা হবে এবং কোন বিধি তা প্রতিরোধ করে, ব্যাখ্যা করো।
\end{minipage}
\end{enumerate}
\end{multicols}

\textbf{৪. ব্যাংকের অসম্পূর্ণ লেনদেন।} A01 হিসাবে 10000 এবং A02 হিসাবে 3000 টাকা আছে। A01 থেকে A02-তে 2000 টাকা পাঠাতে প্রথম হিসাব থেকে টাকা কমানো হলো। দ্বিতীয় হিসাবে যোগ করার আগে ব্যর্থতা ধরা পড়ল; তখনো COMMIT হয়নি। কোনো ফি নেই।
\begin{multicols}{2}
\begin{enumerate}
\item[ক.] ট্রানজেকশন কী?
\item[খ.] Atomicity ও Durability আলাদা কেন?
\item[গ.] ROLLBACK-এর পর দুই Balance নির্ণয় করো; সফল হলে কত হতো?
\item[ঘ.] শুধু Debit স্থায়ী করলে কোন সমস্যা হবে? Consistency বলতে সব লেনদেনে মোট টাকা অপরিবর্তিত বোঝায় কি?
\end{enumerate}
\end{multicols}

\textbf{৫. দোকানের সমন্বিত ব্যবস্থা।} অর্ডার গ্রহণের পর বিক্রয় বিভাগ গ্রাহকের তথ্য, গুদাম বিভাগ পণ্যের পরিমাণ এবং হিসাব বিভাগ পরিশোধের তথ্য ব্যবহার করে। একই গ্রাহকের ঠিকানা তিন বিভাগে আলাদা থাকায় পণ্য ভুল ঠিকানায় গেছে।
\begin{multicols}{2}
\begin{enumerate}
\item[ক.] \begin{minipage}[t]{\linewidth}রিডানড্যান্সি কী?
\end{minipage}
\item[খ.] \begin{minipage}[t]{\linewidth}একই ডেটার ভিন্ন সংস্করণ কেন সমস্যা তৈরি করে?
\end{minipage}
\item[গ.] \begin{minipage}[t]{\linewidth}CustomerID ও OrderID ব্যবহার করে তথ্য সংগঠনের ধারণা দাও।
\end{minipage}
\item[ঘ.] \begin{minipage}[t]{\linewidth}ERP-ধরনের সমন্বয় সুবিধা দিলেও অনুমতি ও যাচাই দরকার কেন, বিশ্লেষণ করো।
\end{minipage}
\end{enumerate}
\end{multicols}

\textbf{৬. ফলাফলের অনুমতি।} শিক্ষার্থী নিজের ফল পড়ে, শিক্ষক নির্দিষ্ট বিষয়ের নম্বর সংশোধন করেন এবং দায়িত্বপ্রাপ্ত কর্মী নতুন শিক্ষার্থীর রেকর্ড যোগ করেন। ভুল কনফিগারেশনে সবার Delete অধিকার চালু হয়ে গেল।
\begin{multicols}{2}
\begin{enumerate}
\item[ক.] \begin{minipage}[t]{\linewidth}Authorization কী?
\end{minipage}
\item[খ.] \begin{minipage}[t]{\linewidth}পরিচয় যাচাই হয়ে গেলেই সব কাজের অধিকার দেওয়া যায় না কেন?
\end{minipage}
\item[গ.] \begin{minipage}[t]{\linewidth}তিন ধরনের ব্যবহারকারীর প্রয়োজনীয় অনুমতির সারণি তৈরি করো।
\end{minipage}
\item[ঘ.] \begin{minipage}[t]{\linewidth}অতিরিক্ত Delete অধিকার, ব্যাকআপ ও পরিবর্তনের লগের সম্পর্ক বিশ্লেষণ করো।
\end{minipage}
\end{enumerate}
\end{multicols}

\textbf{৭. বার্তার রূপান্তর।} একটি শিক্ষামূলক পদ্ধতিতে প্রতিটি ইংরেজি অক্ষর এক ঘর সামনে সরানো হয়; Z-এর পরে A আসে। এ নিয়মে DATA লেখা EBUB হয়েছে। একই প্রতিষ্ঠান বাস্তব ব্যক্তিগত তথ্য সংরক্ষণের নিরাপদ ব্যবস্থা করতে চায়।
\begin{multicols}{2}
\begin{enumerate}
\item[ক.] \begin{minipage}[t]{\linewidth}Ciphertext কী?
\end{minipage}
\item[খ.] \begin{minipage}[t]{\linewidth}এনক্রিপশন ও ডিক্রিপশনের দিক আলাদা কেন?
\end{minipage}
\item[গ.] \begin{minipage}[t]{\linewidth}প্রদত্ত বার্তায় plaintext, ciphertext ও সরণ-কী শনাক্ত করো।
\end{minipage}
\item[ঘ.] \begin{minipage}[t]{\linewidth}এই সরল নিয়মকে বাস্তব ব্যক্তিগত তথ্যের যথেষ্ট নিরাপত্তা বলা যায় কি না, যুক্তি দাও।
\end{minipage}
\end{enumerate}
\end{multicols}

\textbf{৮. অনুসন্ধান ও মেধাতালিকা।} লক্ষাধিক শিক্ষার্থীর তথ্য থেকে কর্মচারী প্রতিদিন RegistrationNo দিয়ে রেকর্ড খোঁজেন। অধ্যক্ষ GPA বড় থেকে ছোট এবং সমান GPA-তে Roll ছোট থেকে বড় তালিকা চান। সব ফিল্ডে ইনডেক্স দেওয়ার প্রস্তাব এসেছে।
\begin{multicols}{2}
\begin{enumerate}
\item[ক.] \begin{minipage}[t]{\linewidth}sort key কী?
\end{minipage}
\item[খ.] \begin{minipage}[t]{\linewidth}ইনডেক্স থাকলেই ফলের ক্রম নিশ্চিত নয় কেন?
\end{minipage}
\item[গ.] \begin{minipage}[t]{\linewidth}অধ্যক্ষের প্রয়োজনের ORDER BY clause লেখো।
\end{minipage}
\item[ঘ.] \begin{minipage}[t]{\linewidth}RegistrationNo-তে ইনডেক্স এবং সব ফিল্ডে ইনডেক্সের প্রস্তাব তুলনা করো।
\end{minipage}
\end{enumerate}
\end{multicols}

\textbf{৯. বিভাগের তথ্য ও নম্বর।} Department(DeptID, DeptName), Student(StudentID, Name, DeptID) এবং Result(ResultID, StudentID, ExamName, ICTMark) আছে। একজন শিক্ষার্থীর একাধিক পরীক্ষার ফল রাখা হবে। একই শিক্ষার্থীর একই পরীক্ষার ফল দ্বিতীয়বার যোগ করা যাবে না।
\begin{multicols}{2}
\begin{enumerate}
\item[ক.] \begin{minipage}[t]{\linewidth}ফরেন কী কী?
\end{minipage}
\item[খ.] \begin{minipage}[t]{\linewidth}Result.StudentID অনন্য করলে সমস্যা কী?
\end{minipage}
\item[গ.] \begin{minipage}[t]{\linewidth}দুটি 1:M সম্পর্ক এবং প্রয়োজনীয় unique বিধি নির্ধারণ করো।
\end{minipage}
\item[ঘ.] \begin{minipage}[t]{\linewidth}Science বিভাগের Annual পরীক্ষায় কমপক্ষে 80 পাওয়া শিক্ষার্থীর নাম ও নম্বর দেখানোর কুয়েরি তৈরি করো।
\end{minipage}
\end{enumerate}
\end{multicols}

\textbf{১০. ব্যাকআপ ও সেবা।} একটি সরকারি সেবা কেন্দ্রে সোমবারের পূর্ণ ব্যাকআপ আছে। মঙ্গলবার নতুন রেকর্ড যোগ হয়েছে। বুধবার সার্ভারের ডিস্ক নষ্ট হলো। সব ডেটা এনক্রিপ্ট করা ছিল, কিন্তু সর্বশেষ পরিবর্তনের আলাদা কপি নেই।
\begin{multicols}{2}
\begin{enumerate}
\item[ক.] \begin{minipage}[t]{\linewidth}Recovery কী?
\end{minipage}
\item[খ.] \begin{minipage}[t]{\linewidth}এনক্রিপশন থাকলেও হারানো ফাইল নিজে থেকে ফিরে আসে না কেন?
\end{minipage}
\item[গ.] \begin{minipage}[t]{\linewidth}বর্তমানে কোন সময় পর্যন্ত তথ্য ফিরে পাওয়ার সম্ভাবনা আছে, ব্যাখ্যা করো।
\end{minipage}
\item[ঘ.] \begin{minipage}[t]{\linewidth}সাম্প্রতিক পরিবর্তন রক্ষা, আলাদা কপি ও পুনরুদ্ধারের পরীক্ষা নিয়ে পরিকল্পনা দাও।
\end{minipage}
\end{enumerate}
\end{multicols}

\subsection{উদ্দীপক থেকে শর্ত ও সংশোধন নির্ণয়}
একটি স্বাস্থ্যকেন্দ্র Patient টেবিলে নিচের তথ্য রেখেছে। কেন্দ্রটি শুধু নির্দিষ্ট রোগীর বয়স সংশোধন এবং নির্দিষ্ট রক্তের গ্রুপের প্রাপ্তবয়স্ক রোগীদের তালিকা তৈরি করতে চায়।
\begin{flushleft}\small
\begin{tabular}{|l|l|l|l|}\hline
PatientID & Name & Age & BloodGroup\\\hline
31 & Rumi & 22 & A+\\\hline
32 & Sajal & 17 & O+\\\hline
33 & Dina & 45 & O+\\\hline
34 & Nafis & 29 & B+\\\hline
\end{tabular}
\end{flushleft}
\begin{enumerate}
\item[ক.] অ্যাট্রিবিউট কী?
\item[খ.] রোগীর নামকে প্রাইমারি কী না করার কারণ ব্যাখ্যা করো।
\item[গ.] Sajal-এর সঠিক বয়স 19 হলে নির্দিষ্ট রেকর্ডটি সংশোধনের SQL লেখো।
\item[ঘ.] সংশোধনের পরে O+ গ্রুপের কমপক্ষে 18 বছর বয়সী রোগীদের বয়সের অধঃক্রমে দেখানোর SQL ও ফল নির্ণয় করো।
\end{enumerate}
\subsection{শেষ অনুশীলন: নতুন টেবিলে নিজে কুয়েরি লেখো}


\subsubsection{অনুশীলনের টেবিল--১: Worker}


\begin{flushleft}

\small

\begin{tabular}{|l|l|l|l|l|l|}
\hline

\chSixHeader
WorkerID & Name & Department & Salary & City & JoinYear\\
\hline
W01 & Rakib & Production & 28000 & Dhaka & 2022\\\hline
W02 & Sumaiya & Accounts & 35000 & Rajshahi & 2021\\\hline
W03 & Imran & Production & 32000 & Khulna & 2023\\\hline
W04 & Tania & Sales & 42000 & Dhaka & 2020\\\hline
W05 & Sakib & Sales & 30000 & Rajshahi & 2022\\\hline
W06 & Mahi & Accounts & 38000 & Dhaka & 2021\\
\hline

\end{tabular}

\end{flushleft}


\begin{chSixSideBySide}[.48]
\begin{enumerate}

\item Worker টেবিলের সব কর্মীর তথ্য দেখাও।

\item শুধু Name ও Salary দেখাও।

\item যেসব কর্মীর City-এর মান Dhaka, তাদের তথ্য দেখাও।

\item যেসব কর্মীর Salary 30000-এর বেশি, তাদের তথ্য দেখাও।

\item যেসব কর্মীর Department-এর মান Production অথবা Sales, তাদের তথ্য দেখাও।

\end{enumerate}
\chSixNextColumn
\begin{enumerate}[start=6]
\item যেসব কর্মীর Salary 28000 থেকে 35000-এর মধ্যে, তাদের তথ্য দেখাও।

\item যেসব কর্মীর JoinYear-এর মান 2022, তাদের তথ্য দেখাও।

\item Salary বড় থেকে ছোট ক্রমে সাজাও।

\item Department-এর মান Accounts এমন কর্মীদের গড় Salary বের করো।

\item সর্বোচ্চ Salary বের করো।

\end{enumerate}
\end{chSixSideBySide}


\subsubsection{অনুশীলনের টেবিল--২: Product}


\begin{flushleft}

\small

\begin{tabular}{|l|l|l|l|l|}
\hline

\chSixHeader
ProductID & ProductName & Category & Price & Stock\\
\hline
P01 & Keyboard & Electronics & 1200 & 15\\\hline
P02 & Mouse & Electronics & 650 & 8\\\hline
P03 & Notebook & Stationery & 120 & 40\\\hline
P04 & Pen Set & Stationery & 250 & 5\\\hline
P05 & Headphone & Electronics & 1800 & 12\\
\hline

\end{tabular}

\end{flushleft}


\begin{multicols}{2}
\begin{enumerate}

\item যেসব পণ্যের Stock 10-এর কম, তাদের তথ্য দেখাও।

\item যেসব পণ্যের Price 500 থেকে 1500-এর মধ্যে, তাদের তথ্য দেখাও।

\item যেসব পণ্যের Category-এর মান Electronics, তাদের তথ্য দেখাও।

\item মোট Stock বের করো।

\item Lowest Price বের করো।

\item Highest Price বের করো।

\item Price Descending Order-এ সাজাও।

\item যেসব পণ্যের ProductName H দিয়ে শুরু, তাদের তথ্য দেখাও।

\end{enumerate}
\end{multicols}


\subsubsection{অনুশীলনের টেবিল--৩: Course}


\begin{flushleft}

\small

\begin{tabular}{|l|l|l|l|}
\hline

\chSixHeader
CourseID & CourseName & Credit & Teacher\\
\hline
C101 & Database & 3 & Salma\\\hline
C102 & Programming & 4 & Rahat\\\hline
C103 & Networking & 3 & Farid\\\hline
C104 & Web Design & 3 & Salma\\\hline
C105 & AI Basics & 4 & Nabila\\
\hline

\end{tabular}

\end{flushleft}


\begin{multicols}{2}
\begin{enumerate}

\item যেসব কোর্সের Credit-এর মান 4, তাদের তথ্য দেখাও।

\item Salma যে কোর্সগুলো পড়ান, সেগুলোর তথ্য দেখাও।

\item যেসব কোর্সের CourseName D দিয়ে শুরু, তাদের তথ্য দেখাও।

\item মোট Course সংখ্যা বের করো।

\item Teacher অনুযায়ী Course সংখ্যা বের করো।

\item CourseName Ascending Order-এ সাজাও।

\end{enumerate}
\end{multicols}


% ============================================================
% PRACTICE HOSPITAL
% ============================================================


\subsubsection{অনুশীলনের টেবিল--৪: Hospital}


\begin{flushleft}

\small

\begin{tabular}{|l|l|l|l|l|}
\hline

\chSixHeader
PatientID & Name & Age & BloodGroup & Bill\\
\hline
H01 & Rafi & 35 & A+ & 12000\\\hline
H02 & Nila & 22 & O+ & 8000\\\hline
H03 & Hasan & 60 & AB+ & 25000\\\hline
H04 & Mitu & 48 & B+ & 18000\\\hline
H05 & Sami & 65 & O+ & 30000\\
\hline

\end{tabular}

\end{flushleft}


\begin{multicols}{2}
\begin{enumerate}

\item যেসব রোগীর Age 50-এর বেশি, তাদের তথ্য দেখাও।

\item যেসব রোগীর রক্তের গ্রুপ O+, তাদের তথ্য দেখাও।

\item যেসব রোগীর Bill 10000 থেকে 25000-এর মধ্যে, তাদের তথ্য দেখাও।

\item Highest Bill বের করো।

\item Average Bill বের করো।

\item Bill Descending Order-এ সাজাও।

\item যেসব রোগীর Name S দিয়ে শুরু, তাদের তথ্য দেখাও।

\end{enumerate}
\end{multicols}


% ============================================================
% FINAL PROBLEM SOLVING MAP
% ============================================================



\subsection{কুয়েরির ধারণা যাচাই}
\begin{practicebox}
\textbf{নিজে যাচাই করো।} কুয়েরি ও SQL অংশের মূল Student ও FeeList সারণি ব্যবহার করো। (১) Student সারণিতে COUNT(*) ও COUNT(Email) কেন আলাদা? (২) FeeList-এ WHERE TuitionFee>1000 এবং HAVING AVG(TuitionFee)>1100 একই ফল বাছাই করে কি?

\end{practicebox}

\subsection{আরও দুটি বহুনির্বাচনি প্রয়োগ}
\begin{enumerate}
\item Science তিনবার, Humanities দুবার এবং Business Studies তিনবার থাকলে \texttt{SELECT DISTINCT Department FROM Student;} কয়টি সারি দেখাবে?\par
ক. 2\quad খ. 3\quad গ. 5\quad ঘ. 8\par
\chSixMCQAnswer{খ}
\item CREATE, ALTER ও DROP-কে কোন শ্রেণিতে রাখা হয়?\par
ক. DDL\quad খ. DCL\quad গ. TCL\quad ঘ. QBE\par
\chSixMCQAnswer{ক}
\end{enumerate}


\subsection{QBE-তে শর্তের অবস্থান: প্রয়োগমূলক MCQ}
\begin{enumerate}
\item \begin{minipage}[t]{\linewidth}Department-এর Criteria সারিতে Science, কিন্তু GPA-এর Or সারিতে \texttt{>=4.80} দেওয়া আছে। কোন রেকর্ড নির্বাচিত হবে?
\begin{tabularx}{\linewidth}{@{}Y Y@{}}
ক. শুধু Science ও GPA অন্তত 4.80 & খ. Science অথবা GPA অন্তত 4.80\\
গ. শুধু Science-এর বাইরের & ঘ. শুধু GPA ঠিক 4.80\\
\end{tabularx}
\chSixMCQAnswer{খ}
\end{minipage}
\item \begin{minipage}[t]{\linewidth}Department-এর Show বন্ধ, কিন্তু Criteria-তে Science আছে। কোনটি সঠিক?
\begin{tabularx}{\linewidth}{@{}Y Y@{}}
ক. শর্ত কাজ করবে না & খ. মূল কলাম মুছে যাবে\\
গ. শর্ত কাজ করবে, কলাম দেখাবে না & ঘ. সব বিভাগ আসবে\\
\end{tabularx}
\chSixMCQAnswer{গ}
\end{minipage}
\end{enumerate}


\subsection{ACID, NoSQL ও Big Data-এর অনুশীলন}
\textbf{ACID শনাক্ত করো।} (১) নিষিদ্ধ ঋণাত্মক Balance গ্রহণ করা হলো না। (২) একই হিসাব থেকে সমসাময়িক দুটি উত্তোলনের সংঘাত সামলানো হলো। (৩) Credit সম্পন্ন না হওয়ায় Debit-ও বাতিল হলো। (৪) সফল COMMIT-এর ফল পুনরায় চালুর পরে পাওয়া গেল।

\textbf{সমন্বিত CQ।} একটি অনলাইন দোকানে বই, পোশাক ও ফোনের বৈশিষ্ট্য আলাদা। প্রতি মুহূর্তে অর্ডার, ছবি ও ক্রেতাদের মতামত আসছে। কিছু মতামত ভুয়া। প্রতিষ্ঠান চাহিদা বুঝে মজুত ঠিক করতে চায়।
\textbf{প্রশ্ন}\par
ক. NoSQL কী?\par
খ. Big Data ও NoSQL একই ধারণা নয় কেন?\par
গ. পণ্যের ভিন্ন বৈশিষ্ট্য রাখার একটি উপযোগী মডেল যুক্তিসহ নির্বাচন করো।\par
ঘ. উদ্দীপকে Velocity, Variety, Veracity ও Value চিহ্নিত করো।


% Topic-specific checks now appear under their canonical theory topics.


Student(StudentID, Name) ও Result(ResultID, StudentID, Marks) থেকে বেশি নম্বর পাওয়া শিক্ষার্থীর নাম দেখাতে কী কী ধারণা লাগবে?

\subsection{সম্পর্ক ও কী: ধারণা যাচাই}
\textbf{নির্বাচিত MCQ}
\begin{chSixSideBySide}[.48]
\small 1. একটি বিভাগের একাধিক শিক্ষার্থী থাকলে বিভাগ থেকে শিক্ষার্থীর সম্পর্ক কোনটি?

\par
\begin{tabularx}{\linewidth}{@{}Y Y@{}}
ক. 1:1 & খ. 1:M \\
গ. M:N & ঘ. 0:1 \\
\end{tabularx}
\smallskip
\chSixMCQAnswer{খ}
\chSixNextColumn
\small 2. Many-to-Many Relationship বাস্তবায়নে সাধারণত কোনটি ব্যবহার করা হয়?

\par
\begin{tabularx}{\linewidth}{@{}Y Y@{}}
ক. Report & খ. Junction Table \\
গ. Index & ঘ. Form \\
\end{tabularx}
\smallskip
\chSixMCQAnswer{খ}
\end{chSixSideBySide}
\begin{chSixSideBySide}[.48]
\small 3. Child Table-এর যে ফিল্ড Parent-এর Key নির্দেশ করে তাকে কী বলে?

\par
\begin{tabularx}{\linewidth}{@{}Y Y@{}}
ক. Primary Key & খ. Foreign Key \\
গ. Sort Key & ঘ. Index Key \\
\end{tabularx}
\smallskip
\chSixMCQAnswer{খ}
\chSixNextColumn
\small 4. Parent Table-এ নেই এমন Foreign Key মান প্রতিরোধের সঙ্গে
কোন ধারণা সবচেয়ে বেশি সম্পর্কিত?

\par
\begin{tabularx}{\linewidth}{@{}Y Y@{}}
ক. Sorting & খ. Referential Integrity \\
গ. Report & ঘ. Form \\
\end{tabularx}
\smallskip
\chSixMCQAnswer{খ}
\end{chSixSideBySide}


\begin{practicebox}

\textbf{নিজে যাচাই করো}

একটি College ডেটাবেজে নিচের টেবিল আছে---

\noindent\texttt{Department(DeptID, DeptName)}\par

\noindent\texttt{Student(StudentID, Name, DeptID)}\par

\noindent\texttt{Result(ResultID, StudentID, ICTMark)}\par

প্রতিটি শিক্ষার্থী একটি Department-এর অন্তর্ভুক্ত। একজন শিক্ষার্থীর
একাধিক পরীক্ষার ফল আলাদা ResultID দিয়ে সংরক্ষিত হতে পারে।

প্রশ্ন:

\begin{enumerate}
    \item কোন ফিল্ডগুলো Primary Key হতে পারে?
    \item কোন ফিল্ডগুলো Foreign Key?
    \item Department ও Student-এর সম্পর্ক কী?
    \item Student ও Result-এর সম্পর্ক কী?
    \item Result থেকে Student-এর সম্পর্ক কোন দিক থেকে কী নামে পরিচিত?
\end{enumerate}

\end{practicebox}





\par\nointerlineskip\vspace{4pt}\noindent\begin{minipage}{\linewidth}
\subsection{সমন্বিত CQ: অনুমতি ও পুনরুদ্ধার}
কলেজের ফলাফলের ডেটাবেজের সর্বশেষ ব্যাকআপ নেওয়া হয়েছিল সোমবার বিকেল ৫টায়। মঙ্গলবার নতুন ফল যোগ করার পর ডিস্ক নষ্ট হয়ে যায়। শিক্ষক নিজের বিষয়ের ফল সংশোধন করতে পারেন; প্রকাশনা কর্মকর্তা শুধু অনুমোদিত ফল পড়তে পারেন। সোমবারের ব্যাকআপ ব্যবহার করে ডেটাবেজটি পুনরুদ্ধার করা হলো।
\textbf{প্রশ্ন}
\begin{enumerate}[label=(\alph*)]
\item ব্যাকআপ কী?
\item সঠিক পাসওয়ার্ড থাকলেও প্রকাশনা কর্মকর্তা ফল বদলাতে পারেন না কেন?
\item ডেটাবেজ ফিরিয়ে আনার পরে কোন তথ্যগুলো যাচাই করতে হবে?
\item সোমবারের ব্যাকআপ-ই মঙ্গলবারের সব তথ্য ফেরানোর নিশ্চয়তা দেয়—বিশ্লেষণ করো।
\end{enumerate}
\end{minipage}\par
% Preserved practice moved from theory to the end-of-chapter bank.

\subsection{Data Security: ধারণা যাচাই}
\begin{multicols}{2}
\small
\begin{enumerate}
\item \begin{minipage}[t]{\linewidth}
একজন শিক্ষার্থী নিজের ফল দেখতে পারে, কিন্তু পরিবর্তন
          করতে পারে না। তাকে কোন ধরনের Authorization দেওয়া হয়েছে?
\end{minipage}
\item \begin{minipage}[t]{\linewidth}
একজন শিক্ষক সঠিক পরিচয়ে লগইন করেছেন, কিন্তু অন্য বিভাগের ফল
          পরিবর্তনের অনুমতি পাননি। এই কাজের সীমা নির্ধারণে Authentication না
          Authorization কাজ করছে?
\end{minipage}
\item \begin{minipage}[t]{\linewidth}
ডিস্ক নষ্ট হয়ে ফলাফলের ডেটাবেজ হারিয়ে গেছে।
          পাসওয়ার্ড জানা থাকলে কি হারানো ডেটা ফিরে পাওয়া যাবে? কেন?
\end{minipage}
\item \begin{minipage}[t]{\linewidth}
ডেটাবেজে GPA-এর রেকর্ড আছে, কিন্তু সংরক্ষিত মান ভুল হয়ে গেছে।
          এটি Data Loss নাকি Data Corruption?
\end{minipage}
\item \begin{minipage}[t]{\linewidth}
কে কখন ফল পরিবর্তন করেছে তা জানতে নিরাপত্তার কোন ব্যবস্থা
          সবচেয়ে প্রাসঙ্গিক?
\end{minipage}
\item \begin{minipage}[t]{\linewidth}
Mobile Number
          \texttt{017*****321}
          আকারে দেখানো হলে কোন নিরাপত্তাপদ্ধতি ব্যবহার করা হয়েছে?
\end{minipage}
\end{enumerate}
\end{multicols}



\textbf{নির্বাচিত MCQ}
\begin{chSixSideBySide}[.48]
\small 1. ব্যবহারকারীর পরিচয় যাচাই করে কোনটি?

\par
\begin{tabularx}{\linewidth}{@{}Y Y@{}}
ক. Authorization & খ. Authentication \\
গ. Recovery & ঘ. Masking \\
\end{tabularx}
\smallskip
\chSixMCQAnswer{খ}
\chSixNextColumn
\small 2. ডেটা শুধু পড়ার অনুমতি কোনটি?

\par
\begin{tabularx}{\linewidth}{@{}Y Y@{}}
ক. Read Authorization & খ. Drop Authorization \\
গ. Alteration Authorization & ঘ. Resource Authorization \\
\end{tabularx}
\smallskip
\chSixMCQAnswer{ক}
\end{chSixSideBySide}
\begin{chSixSideBySide}[.48]
\small 3. সংরক্ষণের ডিস্ক নষ্ট হয়ে ডেটা হারানো কোন ধরনের ব্যর্থতার উদাহরণ?

\par
\begin{tabularx}{\linewidth}{@{}Y Y@{}}
ক. Transaction Failure & খ. Disk Failure \\
গ. Authorization Failure & ঘ. Read Failure \\
\end{tabularx}
\smallskip
\chSixMCQAnswer{খ}
\chSixNextColumn
\small 4. কে কখন রেকর্ড পরিবর্তন করেছে তা জানতে সবচেয়ে উপযোগী কোনটি?

\par
\begin{tabularx}{\linewidth}{@{}Y Y@{}}
ক. Data Masking & খ. Audit Log \\
গ. Drop Authorization & ঘ. Sorting \\
\end{tabularx}
\smallskip
\chSixMCQAnswer{খ}
\end{chSixSideBySide}


\Needspace{7\baselineskip}

\subsection{ডেটা এনক্রিপশন: ধারণা যাচাই}
\begin{multicols}{2}
\small
\begin{enumerate}
\item \begin{minipage}[t]{\linewidth}
Plaintext এবং Ciphertext-এর পার্থক্য কী?
\end{minipage}
\item \begin{minipage}[t]{\linewidth}
Symmetric এনক্রিপশনে একই কী ব্যবহারের প্রধান
          সুবিধা ও সমস্যা কী?
\end{minipage}
\item \begin{minipage}[t]{\linewidth}
Asymmetric পদ্ধতিতে Public Key ও Private Key-এর কাজ কী?
\end{minipage}
\item \begin{minipage}[t]{\linewidth}
Caesar Cipher-এ \(k=3\) হলে \texttt{ABC} কী হবে?
\end{minipage}
\item \begin{minipage}[t]{\linewidth}
এনক্রিপ্ট করা ব্যাকআপ ফাইলের কী হারিয়ে গেলে কী সমস্যা হবে?
\end{minipage}
\item \begin{minipage}[t]{\linewidth}
কেন এনক্রিপশনকে ব্যাকআপের বিকল্প বলা যায় না?
\end{minipage}
\end{enumerate}
\end{multicols}



\textbf{নির্বাচিত MCQ}
\begin{chSixSideBySide}[.48]
\small 1. এনক্রিপশনের আগে মূল বার্তাকে কী বলা হয়?

\par
\begin{tabularx}{\linewidth}{@{}Y Y@{}}
ক. Ciphertext & খ. Plaintext \\
গ. Key & ঘ. Algorithm \\
\end{tabularx}
\smallskip
\chSixMCQAnswer{খ}
\chSixNextColumn
\small 2. একই Secret Key দিয়ে এনক্রিপশন ও ডিক্রিপশন কোন পদ্ধতিতে হয়?

\par
\begin{tabularx}{\linewidth}{@{}Y Y@{}}
ক. Asymmetric & খ. Symmetric \\
গ. Public-only & ঘ. Hashing \\
\end{tabularx}
\smallskip
\chSixMCQAnswer{খ}
\end{chSixSideBySide}
\begin{chSixSideBySide}[.48]
\small 3. কোনটি Asymmetric Algorithm-এর পরিচিত উদাহরণ?

\par
\begin{tabularx}{\linewidth}{@{}Y Y@{}}
ক. AES & খ. DES \\
গ. RSA & ঘ. IDEA \\
\end{tabularx}
\smallskip
\chSixMCQAnswer{গ}
\chSixNextColumn
\small 4. কোন Key গোপন রাখতে হয়?

\par
\begin{tabularx}{\linewidth}{@{}Y Y@{}}
ক. Public Key & খ. Private Key \\
গ. দুটিই প্রকাশ্য & ঘ. কোনোটিই নয় \\
\end{tabularx}
\smallskip
\chSixMCQAnswer{খ}
\end{chSixSideBySide}


\subsection{অতিরিক্ত সংক্ষিপ্ত প্রশ্ন}
\subsubsection{Data, Information ও Database}
\begin{enumerate}
\item একই শিক্ষার্থীর Address তিনটি আলাদা খাতায় রাখা হলো।
    একটি খাতায় Address পরিবর্তন করা হলো, অন্য দুটিতে করা হলো না।
    কী সমস্যা সৃষ্টি হবে?

    \item প্রকৃত নম্বর (Marks) 85-এর পরিবর্তে ডেটাবেজে 58 লেখা হলো।
    পরে ফলাফল (Result) তৈরি করলে কী সমস্যা হতে পারে?

    \end{enumerate}

\subsubsection{Data ও Database: সমন্বিত প্রয়োগ}
\noindent\textbf{CQ ধরনের ধারণামূলক প্রশ্ন।} একটি কলেজে ভর্তি, উপস্থিতি ও ফলাফলের তথ্য তিনটি আলাদা রেজিস্টারে
রাখা হয়। একই শিক্ষার্থীর মোবাইল নম্বর তিন জায়গায় লেখা আছে।
মোবাইল নম্বর পরিবর্তনের পর শুধু ভর্তি রেজিস্টারে নতুন নম্বর লেখা হলো।

\textbf{প্রশ্ন:} এই পদ্ধতির সমস্যা কী এবং ডেটাবেজ ব্যবহার করলে
কীভাবে সমস্যা কমানো যায়?





\subsubsection{Sorting}
\begin{enumerate}
\item Query Result সাজানোর Clause কোনটি?

    \item \texttt{ORDER BY Roll;} কোন Order নির্দেশ করে?

    \item GPA বেশি থেকে কম সাজাতে কী লিখবে?

    \item Text Ascending-এর সাধারণ অর্থ কী?

    \item Date Descending-এর সাধারণ অর্থ কী?

    \item একাধিক ফিল্ডে সর্টিং করলে কোন ফিল্ডের অগ্রাধিকার বেশি?

    \end{enumerate}

